\documentclass[10pt,a4paper]{amsart}
\usepackage{amsmath,amssymb,amsthm,graphicx}
% For commutative diagrams we use Xy-pic
\usepackage[all]{xy}

% We use 2cell for 2-commutative diagrams.
\xyoption{2cell}
% Theorem environments.
%
\theoremstyle{plain}
\newtheorem{theorem}[subsection]{Theorem}
\newtheorem{proposition}[subsection]{Proposition}
\newtheorem{lemma}[subsection]{Lemma}

\theoremstyle{definition}
\newtheorem{definition}[subsection]{Definition}
\newtheorem{example}[subsection]{Example}
\newtheorem{exercise}[subsection]{Exercise}
\newtheorem{situation}[subsection]{Situation}

\theoremstyle{remark}
\newtheorem{remark}[subsection]{Remark}
\newtheorem{remarks}[subsection]{Remarks}

\numberwithin{equation}{subsection}

% Macros
%
\def\lim{\mathop{\mathrm{lim}}\nolimits}
\def\colim{\mathop{\mathrm{colim}}\nolimits}
\def\Spec{\mathop{\mathrm{Spec}}}
\def\Hom{\mathop{\mathrm{Hom}}\nolimits}
\def\Ext{\mathop{\mathrm{Ext}}\nolimits}
\def\SheafHom{\mathop{\mathcal{H}\!\mathit{om}}\nolimits}
\def\SheafExt{\mathop{\mathcal{E}\!\mathit{xt}}\nolimits}
\def\Sch{\mathit{Sch}}
\def\Mor{\mathop{\mathrm{Mor}}\nolimits}
\def\Ob{\mathop{\mathrm{Ob}}\nolimits}
\def\Sh{\mathop{\mathit{Sh}}\nolimits}
\def\NL{\mathop{N\!L}\nolimits}
\def\CH{\mathop{\mathrm{CH}}\nolimits}
\def\proetale{{pro\text{-}\acute{e}tale}}
\def\etale{{\acute{e}tale}}
\def\QCoh{\mathit{QCoh}}
\def\Ker{\mathop{\mathrm{Ker}}}
\def\Im{\mathop{\mathrm{Im}}}
\def\Coker{\mathop{\mathrm{Coker}}}
\def\Coim{\mathop{\mathrm{Coim}}}

% Boxtimes
%
\DeclareMathSymbol{\boxtimes}{\mathbin}{AMSa}{"02}

%
% Macros for moduli stacks/spaces
%
\def\QCohstack{\mathcal{QC}\!\mathit{oh}}
\def\Cohstack{\mathcal{C}\!\mathit{oh}}
\def\Spacesstack{\mathcal{S}\!\mathit{paces}}
\def\Quotfunctor{\mathrm{Quot}}
\def\Hilbfunctor{\mathrm{Hilb}}
\def\Curvesstack{\mathcal{C}\!\mathit{urves}}
\def\Polarizedstack{\mathcal{P}\!\mathit{olarized}}
\def\Complexesstack{\mathcal{C}\!\mathit{omplexes}}
% \Pic is the operator that assigns to X its picard group, usage \Pic(X)
% \Picardstack_{X/B} denotes the Picard stack of X over B
% \Picardfunctor_{X/B} denotes the Picard functor of X over B
\def\Pic{\mathop{\mathrm{Pic}}\nolimits}
\def\Picardstack{\mathcal{P}\!\mathit{ic}}
\def\Picardfunctor{\mathrm{Pic}}
\def\Deformationcategory{\mathcal{D}\!\mathit{ef}}
\usepackage{cleveref}
\crefname{originaltheorem}{original theorem}{original theorems}
\Crefname{originaltheorem}{Original theorem}{Original theorems}
\makeatletter
\newlabel{theorem-picard-group}{{03P8}{source}{}{sourcepointer-1}{}}
\newlabel{theorem-picard-group@cref}{{[originaltheorem][0][]03P8}{[1][0][]source}{}{}{}}
\makeatother
\crefname{originaltheorem}{original theorem}{original theorems}
\Crefname{originaltheorem}{Original theorem}{Original theorems}
\makeatletter
\newlabel{theorem-vanishing-cohomology-Gm-curve}{{03RM}{source}{}{sourcepointer-2}{}}
\newlabel{theorem-vanishing-cohomology-Gm-curve@cref}{{[originaltheorem][0][]03RM}{[1][0][]source}{}{}{}}
\makeatother
\crefname{originallemma}{original lemma}{original lemmas}
\Crefname{originallemma}{Original lemma}{Original lemmas}
\makeatletter
\newlabel{lemma-kummer-sequence}{{03PL}{source}{}{sourcepointer-3}{}}
\newlabel{lemma-kummer-sequence@cref}{{[originallemma][0][]03PL}{[1][0][]source}{}{}{}}
\makeatother
\crefname{originallemma}{original lemma}{original lemmas}
\Crefname{originallemma}{Original lemma}{Original lemmas}
\makeatletter
\newlabel{pic-lemma-picard-pieces}{{0BA0}{source}{}{sourcepointer-4}{}}
\newlabel{pic-lemma-picard-pieces@cref}{{[originallemma][0][]0BA0}{[1][0][]source}{}{}{}}
\makeatother
\crefname{originallemma}{original lemma}{original lemmas}
\Crefname{originallemma}{Original lemma}{Original lemmas}
\makeatletter
\newlabel{groupoids-definition-abelian-variety}{{03RO}{source}{}{sourcepointer-5}{}}
\newlabel{groupoids-definition-abelian-variety@cref}{{[originallemma][0][]03RO}{[1][0][]source}{}{}{}}
\makeatother
\crefname{originallemma}{original lemma}{original lemmas}
\Crefname{originallemma}{Original lemma}{Original lemmas}
\makeatletter
\newlabel{groupoids-proposition-review-abelian-varieties}{{03RP}{source}{}{sourcepointer-6}{}}
\newlabel{groupoids-proposition-review-abelian-varieties@cref}{{[originallemma][0][]03RP}{[1][0][]source}{}{}{}}
\makeatother
\title{Cohomology of the roots of unity on a smooth projective curve}
\author{The Stacks Project Authors}
\date{stacks-project:etale-cohomology.tex:local-snapshot-20261011; GNU FDL 1.2 or later}
\begin{document}
\maketitle
\noindent\textit{Editorial note.} All delivered mathematical text is the author's own, copied from the pinned Stacks Project chapter file etale-cohomology.tex; no statement or proof step was rewritten, reordered or removed.\par\medskip
\noindent\textit{Editorial note.} The proof is the author's complete computation: the sequence is run and the four groups are read off, and the additional paragraph about Pic\textasciicircum{}0(X) completes the noncanonical statement.\par\medskip
\noindent\textit{Editorial note.} Every labelled reference occurring inside the counted spans is carried in assembly.external\_references as a non-load-bearing navigation pointer (same-chapter ones at their definition site, cross-chapter ones at the invocation inside the counted span), so no load-bearing mathematics is replaced by a stub. No bibliography citation occurs inside the counted spans.\par\medskip
\noindent\textit{Editorial note.} The checked chapter list for same-name collisions is the set of stacks chapters staged in this repository: chow, cohomology, divisors, etale-cohomology, flat, more-algebra, more-morphisms, morphisms, spaces-more-morphisms, varieties together with this card's chapter etale-cohomology.\par\medskip
\noindent\textit{Editorial note.} The preamble is the chapter's own native preamble excerpt already used for this source: the Xy-pic setup (native preamble lines 14-18) and the native theorem declarations with the mathematical macros (native preamble lines 154-217).\par\medskip
\noindent\textit{Allowed standard prerequisites (editorial).}\begin{itemize}
\item Kummer sequence and the identification of the first cohomology with the Picard group. Same chapter, the Kummer sequence 0 -> mu\_n -> G\_m -> G\_m -> 0 and the resulting description of H\textasciicircum{}1 with coefficients mu\_n in terms of the Picard group. Supplies the identifications in degrees one and two and the vanishing beyond degree two.
\item Cohomology of the structure sheaf and of the multiplicative group on a curve. Same chapter, the theorems computing the cohomology of the multiplicative group on a smooth projective curve and the vanishing of the cohomology of the structure sheaf. Supplies the input listed at the start of the proof.
\item Structure of the Picard group of a smooth projective curve of genus g. Same chapter together with the classical description: Pic(X) is an extension of Z by the divisible group Pic\textasciicircum{}0(X) of dimension g, whose n-torsion is (Z/n)\textasciicircum{}\{2g\}. Supplies the noncanonical identifications of the second display.
\end{itemize}\medskip
\ifdefined\hypertarget\hypertarget{sourcepointer-1}{}\fi
\noindent\textit{Source reference.} Original theorem 03P8: sources/stacks-etale-cohomology/pinned-original-source.tex, source lines 2498--2498. Navigation pointer invoked at proof line 11166; the counted span contains the pointer only and the cited result is not reproduced inside it, so the pointer is registered as a non-load-bearing navigation stub in the established form for Stacks material. No bibliography citation is involved.\par\medskip
\ifdefined\hypertarget\hypertarget{sourcepointer-2}{}\fi
\noindent\textit{Source reference.} Original theorem 03RM: sources/stacks-etale-cohomology/pinned-original-source.tex, source lines 11074--11074. Navigation pointer invoked at proof line 11167; the counted span contains the pointer only and the cited result is not reproduced inside it, so the pointer is registered as a non-load-bearing navigation stub in the established form for Stacks material. No bibliography citation is involved.\par\medskip
\ifdefined\hypertarget\hypertarget{sourcepointer-3}{}\fi
\noindent\textit{Source reference.} Original lemma 03PL: sources/stacks-etale-cohomology/pinned-original-source.tex, source lines 2842--2842. Navigation pointer invoked at proof line 11171; the counted span contains the pointer only and the cited result is not reproduced inside it, so the pointer is registered as a non-load-bearing navigation stub in the established form for Stacks material. No bibliography citation is involved.\par\medskip
\ifdefined\hypertarget\hypertarget{sourcepointer-4}{}\fi
\noindent\textit{Source reference.} Original lemma 0BA0: sources/stacks-etale-cohomology/pinned-original-source.tex, source lines 11202--11202. Navigation pointer invoked at proof line 11202; the counted span contains the pointer only and the cited result is not reproduced inside it, so the pointer is registered as a non-load-bearing navigation stub in the established form for Stacks material. No bibliography citation is involved.\par\medskip
\ifdefined\hypertarget\hypertarget{sourcepointer-5}{}\fi
\noindent\textit{Source reference.} Original lemma 03RO: sources/stacks-etale-cohomology/pinned-original-source.tex, source lines 11205--11205. Navigation pointer invoked at proof line 11205; the counted span contains the pointer only and the cited result is not reproduced inside it, so the pointer is registered as a non-load-bearing navigation stub in the established form for Stacks material. No bibliography citation is involved.\par\medskip
\ifdefined\hypertarget\hypertarget{sourcepointer-6}{}\fi
\noindent\textit{Source reference.} Original lemma 03RP: sources/stacks-etale-cohomology/pinned-original-source.tex, source lines 11207--11207. Navigation pointer invoked at proof line 11207; the counted span contains the pointer only and the cited result is not reproduced inside it, so the pointer is registered as a non-load-bearing navigation stub in the established form for Stacks material. No bibliography citation is involved.\par\medskip
\begin{lemma}
\label{lemma-cohomology-smooth-projective-curve}
Let $X$ be a smooth projective curve of genus $g$ over an
algebraically closed field $k$ and let $n \geq 1$ be invertible in $k$.
Then there are canonical identifications
$$
H_\etale^q(X, \mu_n) =
\left\{
\begin{matrix}
\mu_n(k) & \text{ if }q = 0, \\
\Pic^0(X)[n] & \text{ if }q = 1, \\
\mathbf{Z}/n\mathbf{Z} & \text{ if }q = 2, \\
0 & \text{ if }q \geq 3.
\end{matrix}
\right.
$$
Since $\mu_n \cong \underline{\mathbf{Z}/n\mathbf{Z}}$, this gives
(noncanonical) identifications
$$
H_\etale^q(X, \underline{\mathbf{Z}/n\mathbf{Z}}) \cong
\left\{
\begin{matrix}
\mathbf{Z}/n\mathbf{Z} & \text{ if }q = 0, \\
(\mathbf{Z}/n\mathbf{Z})^{2g} & \text{ if }q = 1, \\
\mathbf{Z}/n\mathbf{Z} & \text{ if }q = 2, \\
0 & \text{ if }q \geq 3.
\end{matrix}
\right.
$$
\end{lemma}

\begin{proof}
Theorems \ref{theorem-picard-group} and
\ref{theorem-vanishing-cohomology-Gm-curve}
determine the \'etale cohomology of $\mathbf{G}_m$ on $X$
in terms of the Picard group of $X$.
The Kummer sequence $0\to \mu_{n, X} \to \mathbf{G}_{m, X}
\to \mathbf{G}_{m, X}\to 0$ (Lemma \ref{lemma-kummer-sequence})
then gives us the long exact cohomology sequence
$$
\xymatrix{
0 \ar[r] & \mu_n(k) \ar[r] &
k^* \ar[r]^{(\cdot)^n} &
k^* \ar@(rd, ul)[rdllllr] \\
& H_\etale^1(X, \mu_n) \ar[r] &
\Pic(X) \ar[r]^{(\cdot)^n} &
\Pic(X) \ar@(rd, ul)[rdllllr] \\
& H_\etale^2(X, \mu_n) \ar[r] & 0 \ar[r] & 0 \ldots
}
$$
The $n$th power map $k^* \to k^*$ is surjective since $k$ is algebraically
closed. So we need to compute the kernel and cokernel of the map
$n : \Pic(X) \to \Pic(X)$. Consider the commutative diagram with
exact rows
$$
\xymatrix{
0 \ar[r] &
\Pic^0(X) \ar[r] \ar@{>>}[d]^{(\cdot)^n} &
\Pic(X) \ar[r]_-\deg \ar[d]^{(\cdot)^n} &
\mathbf{Z} \ar[r] \ar@{^{(}->}[d]^n & 0 \\
0 \ar[r] &
\Pic^0(X) \ar[r] &
\Pic(X) \ar[r]^-\deg &
\mathbf{Z} \ar[r] & 0
}
$$
The group $\Pic^0(X)$ is the $k$-points of
the group scheme $\underline{\Picardfunctor}^0_{X/k}$, see
Picard Schemes of Curves, Lemma \ref{pic-lemma-picard-pieces}.
The same lemma tells us that $\underline{\Picardfunctor}^0_{X/k}$
is a $g$-dimensional abelian variety over $k$ as defined in
Groupoids, Definition \ref{groupoids-definition-abelian-variety}.
Hence the left vertical map is surjective by
Groupoids, Proposition \ref{groupoids-proposition-review-abelian-varieties}.
Applying the snake lemma gives canonical identifications as stated
in the lemma.

\medskip\noindent
To get the noncanonical identifications of the lemma we need to
show the kernel of $n : \Pic^0(X) \to \Pic^0(X)$
is isomorphic to $(\mathbf{Z}/n\mathbf{Z})^{\oplus 2g}$.
This is also part of Groupoids, Proposition
\ref{groupoids-proposition-review-abelian-varieties}.
\end{proof}

\end{document}
